% Generated draft: requires mathematical and source review.
\documentclass[11pt]{article}
\usepackage[margin=25mm]{geometry}
\usepackage{fontspec}
\setmainfont{Latin Modern Roman}
\usepackage{amsmath,amssymb,mathtools,unicode-math}
\setmathfont{Latin Modern Math}
\usepackage{enumitem,needspace,xurl,hyperref}
\setlength{\parindent}{0pt}
\setlength{\parskip}{5pt}
\setlength{\emergencystretch}{3em}
\newcommand{\N}{\mathbb{N}}
\newcommand{\Z}{\mathbb{Z}}
\newcommand{\Q}{\mathbb{Q}}
\newcommand{\R}{\mathbb{R}}
\newcommand{\C}{\mathbb{C}}
\newcommand{\F}{\mathbb{F}}
\newcommand{\A}{\mathbb{A}}
\newcommand{\PP}{\mathbb{P}}
\newcommand{\E}{\mathbb{E}}
\newcommand{\eps}{\varepsilon}
\newcommand{\dd}{\,\mathrm{d}}
\begin{document}
\section*{UnsolvedMath ID: 30001724}
\textbf{Source problem code:} catalogue entry
\subsection{Definitions and mathematical statement}
The supplied record poses a historiographical question about the meaning of a \emph{numerical table}, rather than a conjecture with a specified mathematical truth condition.

The proposed description is a list of equal-length lists of numbers expressing a correspondence between \(n\) sets, where \(n\in\Z\) and \(n\geq 2\), together with a proposed mathematical dimension of \(n-1\). If ``correspondence'' is interpreted as a relation between sets \(A_1,\ldots,A_n\), it could be represented by
\[
R\subseteq A_1\times\cdots\times A_n.
\]
This relational interpretation is only a possible clarification; the source does not specify whether the correspondence is a relation, a function, or another structure. It also does not identify the numerical domains, specify whether the lists must be finite, or explicitly identify their common length with \(n\).

\textbf{Question.} How should a numerical table be defined? In particular, is the description as a list of equal-length lists of numbers expressing a correspondence between \(n\geq 2\) sets an appropriate definition, and in what precise sense should such a table have mathematical dimension \(n-1\)?

The term ``mathematical dimension'' is not defined in the supplied record. Thus \(n-1\) is part of the proposed description under discussion, not an established dimension formula. No uniquely determined formal proposition or criterion for an adequate definition is supplied.
\subsection{Short English statement}
The question asks for a suitable historical and mathematical definition of a numerical table. It specifically asks whether equal-length lists of numbers describing a correspondence between \(n\geq 2\) sets should constitute such a table, with dimension \(n-1\). The meanings of correspondence and dimension remain unspecified.
\subsection{Sources}
\begin{itemize}
\item Supplied local UnsolvedMath registry record, permanent ID \(30001724\), title \emph{Mathematical Definition of Numerical Tables}, exact source problem code \texttt{catalogue entry}, in the Oberwolfach Reports Open Problems collection. The proposed year is 2011; no individual proposer is identified.
\item \emph{Mini-Workshop: History of Numerical and Graphical Tables} (2011), Oberwolfach Reports, pp.\ 639--689, DOI \texttt{10.4171/owr/2011/12}. Supplied source citation and literature-source entry, labeled ``Oberwolfach Report 12/2011'': \url{https://doi.org/10.4171/owr/2011/12}.
\item The record supplies statement status \texttt{corrected\_verified} and reports that the official report confirms this historiographical research question and that an extraction label was removed. These are supplied verification claims, not independent verification performed for this draft.
\item Literature assessment dated \texttt{2026-08-21}: the supplied curation labels the question open and describes it as part of an ongoing historiographical research project, with no single settled definition closing the extracted question. The listed source checked is Oberwolfach Report 12/2011 at the DOI above. This is supplied dated curation, not an independently verified literature assessment; time-sensitive status and cited claims require independent checking.
\end{itemize}
\end{document}
